\documentclass{article}
\usepackage{lmodern}
\usepackage{amsmath}
\usepackage{listings}
\usepackage{fullpage}
\usepackage{parskip}
\usepackage{xcolor}
\lstset{
	basicstyle=\ttfamily,
	columns=fullflexible,
	frame=single,
	breaklines=true,
	postbreak=\mbox{\textcolor{red}{$\hookrightarrow$}\space},
}
\begin{document}
\title{Experiment `p\_4\_1\_c\_powers\_of\_two\_seq' Results}
\author{\today}
\date{}
\maketitle
\textbf{Experiment outcome: }SUCCESS
\\\textbf{Bad responses: }0
\\\textbf{Responses containing}~\texttt{assume}~\textbf{: }0
\\\textbf{Responses containing}~\texttt{expect}~\textbf{: }0
\\\textbf{Resolution attempts: }2
\\\textbf{Hard fails (resolution): }1
\\\textbf{Soft fails (resolution): }0
\\\textbf{Verification attempts: }1
\section*{Problem Specification}
\textbf{Problem name: }p\_4\_1\_c\_powers\_of\_two\_seq
\\\textbf{Natural language statement: }Write a method with loops that computes all powers of 2 from 2\^{}0 up to 2\^{}20.
\\\textbf{Method signature: }p\_4\_1\_c\_powers\_of\_two\_seq() returns (powers: seq<int>)
\subsection*{Ensures}
\begin{itemize}
\item \texttt{|powers| == 21}
\item \texttt{forall i :: 0 <= i < |powers| ==> powers[i] == pow(2, i)}
\end{itemize}
\subsection*{Functional Code Given}
\begin{lstlisting}
function pow(base: int, exp: int): int
  requires  0 <= exp
  decreases exp
{
  if exp == 0 then 1 else base * pow(base, exp - 1)
}
\end{lstlisting}
\clearpage
\section*{GenAI interactions}
Below you will find all interactions between the `user' (program) and the `assistant' (OpenAI).
\subsection*{Program $\rightarrow$ GenAI}
\begin{lstlisting}
You are given the following task to perform in Dafny:

Write a method with loops that computes all powers of 2 from 2^0 up to 2^20.

The signature should be:

p_4_1_c_powers_of_two_seq() returns (powers: seq<int>)

The method should respect the following contract:

ensures |powers| == 21, ensures forall i :: 0 <= i < |powers| ==> powers[i] == pow(2, i)

The contract uses the following dafny code:

function pow(base: int, exp: int): int
  requires  0 <= exp
  decreases exp
{
  if exp == 0 then 1 else base * pow(base, exp - 1)
}

These function(s) must not be used in your implementation of the method.

Produce and show only the Dafny body of this method, including the curly braces that surround it. Do not show the signature nor contract. You must not use 'assume' nor 'expect' anywhere in your code.
\end{lstlisting}
\subsection*{GenAI $\rightarrow$ Program}
\textbf{System fingerprint: }null
\\\textbf{ID: }chatcmpl-DmdzQm4Gm323ZyAntEVP5DZFONwpe
\\\textbf{Model: }gpt-5.5-2026-04-23
\\\textbf{Created at: }1780485836
\\\textbf{Finish reason: }stop
\\\textbf{Usage: }468 tokens in, and 6615 tokens out
\begin{lstlisting}
{
  powers := [];
  var i := 0;
  var current := 1;

  while i < 21
    invariant 0 <= i <= 21
    invariant |powers| == i
    invariant current == pow(2, i)
    invariant forall j :: 0 <= j < i ==> powers[j] == pow(2, j)
    decreases 21 - i
  {
    var nextI := i + 1;
    ghost var oldPowers := powers;

    powers := oldPowers + [current];
    assert |powers| == nextI;
    assert powers[i] == current;

    forall j | 0 <= j < nextI
      ensures powers[j] == pow(2, j)
    {
      if j < i {
        assert powers[j] == oldPowers[j];
      } else {
        assert j == i;
        assert powers[j] == current;
      }
    }

    current := current * 2;
    assert current == pow(2, i) * 2;
    assert current == 2 * pow(2, i);
    assert 0 < nextI;
    assert nextI - 1 == i;
    assert pow(2, nextI) == 2 * pow(2, nextI - 1);
    assert pow(2, nextI) == 2 * pow(2, i);
    assert current == pow(2, nextI);

    i := nextI;
  }

  assert i == 21;
  assert |powers| == 21;
  forall j | 0 <= j < |powers|
    ensures powers[j] == pow(2, j)
  {
    assert j < i;
  }
}
\end{lstlisting}
\subsection*{Program $\rightarrow$ GenAI}
\begin{lstlisting}
When using dafny resolve, the below error is emitted and resolve fails:

problem_data/p_4_1_c_powers_of_two_seq/program.dfy(26,14): Error: ghost variables such as oldPowers are allowed only in specification contexts. oldPowers was inferred to be ghost based on its declaration or initialization.
   |
26 |     powers := oldPowers + [current];
   |               ^^^^^^^^^

1 resolution/type errors detected in program.dfy

Correct the error in code by altering only the method body. Produce and show only the Dafny body, including the curly braces that surround it. Do not show the signature nor contract. You must not use 'assume' nor 'expect' anywhere in your code.
\end{lstlisting}
\subsection*{GenAI $\rightarrow$ Program}
\textbf{System fingerprint: }null
\\\textbf{ID: }chatcmpl-Dme1TAGhRbg9B2VUwXpiiABStdXRX
\\\textbf{Model: }gpt-5.5-2026-04-23
\\\textbf{Created at: }1780485963
\\\textbf{Finish reason: }stop
\\\textbf{Usage: }980 tokens in, and 890 tokens out
\begin{lstlisting}
{
  powers := [];
  var i := 0;
  var current := 1;

  while i < 21
    invariant 0 <= i <= 21
    invariant |powers| == i
    invariant current == pow(2, i)
    invariant forall j :: 0 <= j < i ==> powers[j] == pow(2, j)
    decreases 21 - i
  {
    var nextI := i + 1;
    var oldPowers := powers;

    powers := oldPowers + [current];
    assert |powers| == nextI;
    assert powers[i] == current;

    forall j | 0 <= j < nextI
      ensures powers[j] == pow(2, j)
    {
      if j < i {
        assert powers[j] == oldPowers[j];
      } else {
        assert j == i;
        assert powers[j] == current;
      }
    }

    current := current * 2;
    assert current == pow(2, i) * 2;
    assert current == 2 * pow(2, i);
    assert 0 < nextI;
    assert nextI - 1 == i;
    assert pow(2, nextI) == 2 * pow(2, nextI - 1);
    assert pow(2, nextI) == 2 * pow(2, i);
    assert current == pow(2, nextI);

    i := nextI;
  }

  assert i == 21;
  assert |powers| == 21;
  forall j | 0 <= j < |powers|
    ensures powers[j] == pow(2, j)
  {
    assert j < i;
  }
}
\end{lstlisting}
\section*{Final Program}
\begin{lstlisting}
function pow(base: int, exp: int): int
  requires  0 <= exp
  decreases exp
{
  if exp == 0 then 1 else base * pow(base, exp - 1)
}

method p_4_1_c_powers_of_two_seq() returns (powers: seq<int>)
	ensures |powers| == 21
	ensures forall i :: 0 <= i < |powers| ==> powers[i] == pow(2, i)
{
  powers := [];
  var i := 0;
  var current := 1;

  while i < 21
    invariant 0 <= i <= 21
    invariant |powers| == i
    invariant current == pow(2, i)
    invariant forall j :: 0 <= j < i ==> powers[j] == pow(2, j)
    decreases 21 - i
  {
    var nextI := i + 1;
    var oldPowers := powers;

    powers := oldPowers + [current];
    assert |powers| == nextI;
    assert powers[i] == current;

    forall j | 0 <= j < nextI
      ensures powers[j] == pow(2, j)
    {
      if j < i {
        assert powers[j] == oldPowers[j];
      } else {
        assert j == i;
        assert powers[j] == current;
      }
    }

    current := current * 2;
    assert current == pow(2, i) * 2;
    assert current == 2 * pow(2, i);
    assert 0 < nextI;
    assert nextI - 1 == i;
    assert pow(2, nextI) == 2 * pow(2, nextI - 1);
    assert pow(2, nextI) == 2 * pow(2, i);
    assert current == pow(2, nextI);

    i := nextI;
  }

  assert i == 21;
  assert |powers| == 21;
  forall j | 0 <= j < |powers|
    ensures powers[j] == pow(2, j)
  {
    assert j < i;
  }
}
\end{lstlisting}
\section*{Total Token Usage}
\textbf{Input tokens: }1448
\\\textbf{Output tokens: }7505
\\\textbf{Reasoning tokens: }6590
\\\textbf{Sum of `total tokens': }8953
\section*{Experiment Timings}
\textbf{Overall Experiment} started at 1780485835810, ended at 1780485998124, lasting 162314ms (162.31 seconds)
\\\textbf{Iteration \#1} started at 1780485835811, ended at 1780485962951, lasting 127140ms (127.14 seconds)
\\\textbf{Iteration \#2} started at 1780485962956, ended at 1780485998123, lasting 35167ms (35.17 seconds)
\end{document}
